\documentclass[11pt]{article}
\usepackage[a4paper,margin=24mm]{geometry}
\usepackage{amsmath,amssymb,xcolor}
\usepackage[colorlinks=true,linkcolor=blue!45!black,urlcolor=blue!45!black]{hyperref}
\setlength{\parindent}{0pt}
\setlength{\parskip}{6pt}
\setlength{\emergencystretch}{2em}
\newcommand{\Q}{\mathbb Q}
\newcommand{\Z}{\mathbb Z}
\newcommand{\R}{\mathbb R}
\newcommand{\dd}{\,d}
\newcommand{\Hfive}[1]{H_{#1}^{(5)}}
\newcommand{\source}[1]{{\small\color{blue!45!black}#1}\par}
\newcommand{\heading}[2]{\section*{#1}\source{#2}}
\pagestyle{plain}
\begin{document}
\begin{center}
{\LARGE The determinant argument for\\[3pt]
the irrationality of $\zeta(5)$}\\[9pt]
{\large A detailed exposition following Chun Wing Lai}\\[5pt]
Draft \enspace$\cdot$\enspace 14 November 2026
\end{center}

The argument constructs integer polynomials whose values at $\zeta(5)$
become extraordinarily small. Their degrees grow only linearly, while
their values shrink exponentially in the square of the size parameter.
If $\zeta(5)$ were rational, clearing its fixed denominator would produce
a positive integer tending to zero.

This account supplies the details of the construction and that contradiction.
The substantial arithmetic and energy estimates are stated explicitly and
taken as inputs. As in the shorter overview, the arithmetic input is
Fauzan's original bound, assumed separately from Lai's sections 8--11 as
literally printed; the source note explains the distinction. The classical
identities used in the integral calculation are proved in Appendix A.
Appendices B and C give the exact-degree and pullback arguments, which
clarify the construction but are not needed once the two estimates are given.

\heading{1. The contradiction we are aiming for}{Lai, section 3, Lemma 3.1.}
Write $\xi=\zeta(5)$. It is enough to find, for every sufficiently large
positive integer $n$, a polynomial $Q_n\in\Z[X]$ such that
\[
 \deg Q_n\leq37n,\qquad 0<Q_n(\xi)\leq e^{-cn^2},
 \qquad c>0\text{ fixed}.
\]
To see why, suppose $\xi=a/b$, with $a\in\Z$ and $b\in\Z_{>0}$.
Write $Q_n(X)=\sum_{j=0}^{37n}q_jX^j$, adding zero coefficients if necessary.
Then
\[
 b^{37n}Q_n(a/b)=\sum_{j=0}^{37n}q_j a^j b^{37n-j}\in\Z.
\]
This integer is positive. On the other hand,
\[
 0<b^{37n}Q_n(a/b)
 \leq\exp\bigl(37n\log b-cn^2\bigr)\longrightarrow0.
\]
Here $b$ is fixed: the quadratic term eventually dominates the linear term.
A positive integer cannot tend to zero, so the assumed rational expression
for $\xi$ is impossible. The fraction need not be in lowest terms.

The construction must therefore keep a polynomial value both nonzero and
small while making every coefficient integral. Clearing denominators alone
is easy, but a large multiplier can destroy smallness. The determinant is
designed to provide enough decay to pay that arithmetic cost.

\newpage
\heading{2. Defining the functional carefully}{Lai, section 4, Definition 4.2.}
The notation involves three different variables. The letter $t$ is the
argument of a rational function. The letter $X$ is a parameter that will
eventually be set equal to $\xi$. The letter $y$ will be an integration
variable. Keeping these roles separate makes the later matrix calculation
much easier to follow.

Let $\mathcal R$ be the vector space over $\Q$ of rational functions in
$\Q(t)$ whose finite poles, if any, are simple and belong to
$\{-1^2,-2^2,-3^2,\ldots\}$. Every $R\in\mathcal R$ has a unique expression
\begin{equation}\label{eq:partial}
 R(t)=\sum_{e=0}^{d}p_et^e+\sum_{j\in J}\frac{c_j}{t+j^2},
 \qquad p_e,c_j\in\Q,
\end{equation}
where $J$ is finite. Polynomial division gives the polynomial part.
At a pole, its coefficient is
$c_j=\left.(t+j^2)R(t)\right|_{t=-j^2}$.
Subtracting these simple fractions leaves a rational function with no finite
poles, hence a polynomial. These observations also prove uniqueness
(with zero terms omitted).

Define the Bernoulli numbers by
\[
 \frac{u}{e^u-1}=\sum_{k\geq0}B_k\frac{u^k}{k!}.
\]
Multiplication by the exponential series determines the $B_k$ recursively
as rational numbers. In particular $B_0=1$, $B_1=-1/2$, $B_2=1/6$,
and $B_4=-1/30$. Also put
\[
 \Hfive{j}=\sum_{v=1}^j v^{-5},\qquad \Hfive{0}=0.
\]
Prescribe a $\Q$-linear functional $\mu_X:\mathcal R\to\Q[X]$ by
\begin{align}
 \mu_X(t^e)&=
 \frac{(-1)^eB_{2e+2}(2e+3)(2e+4)(2e+5)}{24},
 &&e\in\Z_{\geq0},\label{eq:mu-moment}\\
 \mu_X\!\left(\frac1{t+j^2}\right)&=
 j^4\bigl(X-\Hfive{j}\bigr)-\frac14+\frac1{2j},
 &&j\in\Z_{>0}.\label{eq:mu-pole}
\end{align}
For a general $R$, apply these formulas to its decomposition and add with
the coefficients $p_e,c_j$. This is the meaning of extending by linearity.
It gives a well-defined affine polynomial in $X$: only the simple
fractions contribute a term involving $X$.

For example,
\[
 \mu_X(1)=\frac5{12},\qquad
 \mu_X\!\left(\frac1{t+1}\right)=X-\frac34.
\]
Since $t/(t+1)=1-1/(t+1)$, linearity gives
$\mu_X(t/(t+1))=7/6-X$. There is no rule allowing us to multiply
functional values; $\mu_X$ is linear, not multiplicative.

Finally, $\mu_\xi(R)$ means evaluate the polynomial $\mu_X(R)$ at $X=\xi$.
It does not mean substitute $\xi$ for $t$ in the rational function.

\newpage
\heading{3. Why the polynomial values are moments}{Lai, section 4, Lemma 4.5.}
Define, for $y>0$,
\begin{equation}\label{eq:weight}
 w(y)=\frac{(2\pi)^4y^5}{12}
       \sum_{\ell\geq1}\ell^4e^{-2\pi\ell y}.
\end{equation}
Every summand is positive. The series converges for each $y>0$, so
$w(y)>0$. We will prove
\begin{equation}\label{eq:representation}
 \boxed{\mu_\xi(R)=\int_0^\infty R(y^2)w(y)\dd y
        \quad(R\in\mathcal R).}
\end{equation}
This is the bridge between rational algebra and real integration.
First we check that the integrals make sense.

Put $f(y)=(e^{2\pi y}-1)^{-1}$. Its geometric series may be differentiated
term by term on $y\geq\varepsilon>0$, where the derivative series converge
uniformly. Four differentiations show that
\[
 w(y)=\frac{y^5}{12}f^{(4)}(y).
\]
Near zero, $f(y)=1/(2\pi y)+g(y)$ with $g$ analytic at zero. Consequently
$f^{(4)}(y)=24/(2\pi y^5)+O(1)$ and $w(y)\to1/\pi$.
For $y\geq1$, factor $e^{-2\pi y}$ out of the series in \eqref{eq:weight};
the remaining sum is bounded by
$\sum_{\ell\geq1}\ell^4e^{-2\pi(\ell-1)}<\infty$.
Thus $w(y)=O(y^5e^{-2\pi y})$ at infinity.

An admissible $R(y^2)$ is bounded near zero, because its poles lie away from
$[0,\infty)$. It grows at most polynomially as $y\to\infty$.
These two bounds on $w$ prove absolute convergence in
\eqref{eq:representation}.

For $R(t)=t^e$, the summands of the integral are nonnegative. Tonelli's
theorem allows us to interchange the sum and integral:
\begin{align*}
 \int_0^\infty y^{2e}w(y)\dd y
 &=\frac{(2\pi)^4}{12}\sum_{\ell\geq1}\ell^4
      \int_0^\infty y^{2e+5}e^{-2\pi\ell y}\dd y\\
 &=\frac{(2e+5)!}{12(2\pi)^{2e+2}}
      \sum_{\ell\geq1}\frac1{\ell^{2e+2}}.
\end{align*}
Here $\int_0^\infty y^m e^{-by}\dd y=m!/b^{m+1}$ for $b>0$:
substitution reduces this to $b=1$, and repeated integration by parts gives
$m!$. Euler's formula, proved in Appendix A, is
\[
 \zeta(2e+2)=
 \frac{(-1)^eB_{2e+2}(2\pi)^{2e+2}}{2(2e+2)!}.
\]
Substituting it above yields
\[
 \int_0^\infty y^{2e}w(y)\dd y
 =\frac{(-1)^eB_{2e+2}}{24}\,
       \frac{(2e+5)!}{(2e+2)!}
 =\mu_X(t^e).
\]
The factorial quotient is exactly $(2e+3)(2e+4)(2e+5)$.
The equality is independent of $X$, as the definition requires.

\newpage
\heading{4. Why the pole values are the same integrals}{Lai, section 4, Lemma 4.6.}
The simple fractions need a different calculation. We do it for a real
number $a>0$, then specialize to a positive integer. Define
\[
 q_a(y)=\frac{y^5}{y^2+a^2}
       =y^3-a^2y+a^4\frac{y}{y^2+a^2}.
\]
The polynomial terms disappear after four derivatives. Since
$y/(y^2+a^2)=\operatorname{Im}(a-iy)^{-1}$, and
$\frac{d^4}{dy^4}(a-iy)^{-1}=24(a-iy)^{-5}$, we obtain
\[
 q_a^{(4)}(y)=24a^4\operatorname{Im}(a-iy)^{-5}.
\]
Using $w=y^5f^{(4)}/12$, four integrations by parts give
\begin{align*}
 \int_0^\infty\frac{w(y)}{y^2+a^2}\dd y
 &=\frac1{12}\int_0^\infty q_a(y)f^{(4)}(y)\dd y\\
 &=\frac1{12}\int_0^\infty q_a^{(4)}(y)f(y)\dd y\\
 &=2a^4\int_0^\infty
     \frac{\operatorname{Im}(a-iy)^{-5}}{e^{2\pi y}-1}\dd y.
\end{align*}
There are no missing boundary terms. First work on $[\varepsilon,L]$;
the boundary expression is
$q_af'''-q_a'f''+q_a''f'-q_a'''f$.
For $0\leq r\leq3$, at zero we have
\[
 q_a^{(r)}(y)=O(y^{5-r}),\qquad
 f^{(3-r)}(y)=O(y^{r-4}),
\]
so every boundary product is $O(y)$ and tends to zero. At infinity the
derivatives of $f$ decay exponentially and the other factors have at most
polynomial growth. Also $q_a^{(4)}(y)=O(y)$ at zero, so
$q_a^{(4)}(y)f(y)=O(1)$ there. The final integral converges at both ends,
and we may let $\varepsilon\downarrow0$ and $L\to\infty$.

For $a>0$, define $\zeta(5,a)=\sum_{v=0}^\infty(a+v)^{-5}$.
The particular Hurwitz identity needed here is
\begin{equation}\label{eq:hermite}
 \zeta(5,a)=\frac1{2a^5}+\frac1{4a^4}
 +2\int_0^\infty
     \frac{\operatorname{Im}(a-iy)^{-5}}{e^{2\pi y}-1}\dd y.
\end{equation}
Appendix A derives this directly from an elementary kernel identity.
It gives
\[
 \int_0^\infty\frac{w(y)}{y^2+a^2}\dd y
 =a^4\zeta(5,a)-\frac1{2a}-\frac14.
\]
For an integer $j\geq1$, the series gives
$\zeta(5,j)=\xi-\Hfive{j}+j^{-5}$. Therefore
\[
 \int_0^\infty\frac{w(y)}{y^2+j^2}\dd y
 =j^4(\xi-\Hfive{j})-\frac14+\frac1{2j}
 =\mu_\xi\!\left(\frac1{t+j^2}\right).
\]
We have checked both types of term in \eqref{eq:partial}.
Applying the formulas to its finitely many terms and using linearity
of integration proves \eqref{eq:representation} for every $R\in\mathcal R$.

\newpage
\heading{5. Building a polynomial from the functional}{Lai, section 5, Definitions 5.1--5.2 and Lemma 5.4.}
For a positive integer $n$, choose
\[
 K=40n,\qquad N=3n,\qquad h=K-N=37n,
 \qquad D_m(t)=\prod_{j=1}^m(t+j^2).
\]
The parameters have been chosen to make the later estimates work.
For now, the features that matter are positivity and cancellation. Define
\[
 A_K(t)=\frac{D_N(t)^6}{D_K(t)}
       =\frac{D_N(t)^5}{T_K(t)},\qquad
 T_K(t)=\prod_{k=N+1}^K(t+k^2).
\]
The cancellation removes the first $N$ possible poles. The $h$ remaining
roots of $T_K$ are distinct negative integer squares, so multiplying
$A_K$ by any polynomial leaves a function in $\mathcal R$.
Moreover, $A_K(y^2)>0$ for every $y>0$.

For $0\leq i,j<h$ put $R_{ij}(t)=A_K(t)t^{i+j}$ and define
\begin{equation}\label{eq:G}
 G_K(X)=\bigl[\mu_X(R_{ij})\bigr]_{0\leq i,j<h},
 \qquad \Delta_K(X)=\det G_K(X).
\end{equation}
The dependence on $i+j$ makes this a Hankel matrix: its entries are constant
along each anti-diagonal. The first few entries have the pattern
\[
 \begin{pmatrix}
 \mu_X(A_K)&\mu_X(tA_K)&\mu_X(t^2A_K)&\cdots\\
 \mu_X(tA_K)&\mu_X(t^2A_K)&\mu_X(t^3A_K)&\cdots\\
 \mu_X(t^2A_K)&\mu_X(t^3A_K)&\mu_X(t^4A_K)&\cdots\\
 \vdots&\vdots&\vdots&\ddots
 \end{pmatrix}.
\]
To see explicitly why these are affine entries, take partial fractions:
\[
 R_{ij}(t)=P_{ij}(t)+\sum_{k=N+1}^K\frac{c_k^{ij}}{t+k^2},
 \qquad
 c_k^{ij}=\frac{D_N(-k^2)^5(-k^2)^{i+j}}{T_K'(-k^2)}.
\]
The polynomial part may be zero; it must not otherwise be discarded.
Applying the functional gives
\[
 G_K(X)_{ij}=\mu_X(P_{ij})+
 \sum_{k=N+1}^K c_k^{ij}
 \left[k^4(X-\Hfive{k})-\frac14+\frac1{2k}\right].
\]
Thus every entry is $a_{ij}X+b_{ij}$ with rational $a_{ij},b_{ij}$.
The determinant formula is
\[
 \Delta_K(X)=\sum_{\sigma\in S_h}\operatorname{sgn}(\sigma)
       \prod_{i=0}^{h-1}(a_{i,\sigma(i)}X+b_{i,\sigma(i)}).
\]
Here $S_h$ permutes the indices $0,\ldots,h-1$.
Each product has degree at most $h$, and a finite sum cannot increase that
bound. Hence $\Delta_K\in\Q[X]$ and $\deg\Delta_K\leq37n$.
Appendix B proves exact degree $h$, but the upper bound suffices here.

\newpage
\heading{6. Evaluating the matrix at $\xi$}{Lai, section 5, Proposition 5.6.}
We can now pass from a symbolic matrix to a matrix of integrals one entry
at a time. With $R_{ij}$ as above,
\begin{align*}
 G_K(\xi)_{ij}
 &=\left.\mu_X(R_{ij})\right|_{X=\xi}\\
 &=\mu_\xi(R_{ij})\\
 &=\int_0^\infty R_{ij}(y^2)w(y)\dd y\\
 &=\int_0^\infty y^{2i}y^{2j}A_K(y^2)w(y)\dd y.
\end{align*}
The third equality is the theorem proved in sections 3--4 of this account.
The last is substitution: $(y^2)^{i+j}=y^{2i}y^{2j}$.
Absolute convergence was already checked for every admissible rational
function, including these particular $R_{ij}$.

For a compact description, set
\[
 \dd\nu_K(y)=A_K(y^2)w(y)\dd y,\qquad
 v(y)=(1,y^2,\ldots,y^{2h-2})^{\mathsf T}.
\]
Then $G_K(\xi)=\int_0^\infty v(y)v(y)^{\mathsf T}\dd\nu_K(y)$,
where the integral is taken \emph{entry by entry}. Its $(i,j)$ entry is
exactly the last integral above.

Let $c=(c_0,\ldots,c_{h-1})^{\mathsf T}\in\R^h$ and
$q(t)=\sum_i c_it^i$. Distributing the finite sums gives
\begin{align*}
 c^{\mathsf T}G_K(\xi)c
 &=\sum_{i,j}c_ic_j\int_0^\infty y^{2i+2j}\dd\nu_K(y)\\
 &=\int_0^\infty\left(\sum_{i,j}c_ic_jy^{2i+2j}\right)\dd\nu_K(y)\\
 &=\int_0^\infty\left(\sum_i c_iy^{2i}\right)^2\dd\nu_K(y)
 =\int_0^\infty q(y^2)^2\dd\nu_K(y).
\end{align*}
Only finite linearity is used here. Real coefficients $c_i$ cause no
difficulty with the originally $\Q$-linear definition of $\mu_X$.

If $c\ne0$, then $q$ is a nonzero polynomial. Choose $y_0>0$ that is not
a root of $q(y^2)$. By continuity, $q(y^2)^2$ is bounded below by a positive
number on a small interval about $y_0$. The density of $\nu_K$ is positive
and continuous there, so the integral is strictly positive.

Thus $G_K(\xi)$ is positive definite. To see why its determinant is positive,
orthogonalize $1,y^2,\ldots,y^{2h-2}$ by Gram--Schmidt without normalizing.
The resulting orthogonal functions are monic polynomials in $y^2$, so each
is nonzero and has positive squared norm. The change of basis is triangular
with diagonal entries one. The determinant is therefore the product of
these positive squared norms: $\Delta_K(\xi)>0$.

The matrix integral is not a rule for moving a determinant inside an
integral. Each $v(y)v(y)^{\mathsf T}$ has rank one. The multiple integral
for the determinant comes from a separate identity, stated in section 8 below.

\newpage
\heading{7. Making the coefficients integral at a controlled cost}
{The arithmetic output of Lai's sections 8--10, with the source qualification below.}
The estimates use a fixed explicit rescaling of the determinant:
\[
 F_K(X)=S_K\Delta_K(X),\qquad
 S_K=\frac{(K!)^{2h}4^{h-1}}
 {(N!)^{12h}\displaystyle\prod_{i=1}^{h-1}((2i)!)^2}>0.
\]
All factors are nonzero integers, so $S_K$ is a positive rational number
independent of $X$. Therefore
\[
 F_K\in\Q[X],\qquad \deg F_K\leq h,\qquad F_K(\xi)>0.
\]
The factorial rescaling puts the arithmetic and analytic bounds on a common
scale. It does not by itself assert that all coefficients are integers.

\textbf{Arithmetic input.} Fix $M=200$. There are positive rational
numbers $m_{K,200}$, for all sufficiently large $K\in40\Z_{>0}$, such that
\begin{equation}\label{eq:arith-input}
 m_{K,200}F_K\in\Z[X],\qquad
 \limsup_{\substack{K\to\infty\\40\mid K}}
 \frac{\log m_{K,200}}{K^2}\leq A_{200},
\end{equation}
where
\[
 A_{200}=\frac{127125602969131786927559}
                 {94195881588024216000000}
        =1.349587697741\ldots<1.35.
\]
We assume this arithmetic statement in the form of Fauzan's Proposition 5.1
and equation (5.21), with Appendix B's evaluation. It is not the literal
corrected statement of Lai's Theorem 9.3; the final source note explains why.
The original restrictions $40\mid M$ and $K\geq200M^2$ hold for $M=200$
and sufficiently large $K$. The cutoff stays fixed as $K$ grows.

Here is what the arithmetic assertion accomplishes. For a nonzero rational
number, $v_p$ counts the power of $p$ in the numerator minus its power in
the denominator. A rational number is an integer exactly when all these
exponents are nonnegative. If $L_p$ is an integer lower bound for $v_p$ of
every nonzero coefficient of $F_K$, then multiplying by $p^{-L_p}$ removes
any denominator at $p$; when $L_p>0$, it instead divides out a common factor.
The local estimates justify these bounds and their use for all primes.
Their sum controls the size of the total positive rational multiplier.
We do not re-prove those local estimates here.

Abbreviate $m_K=m_{K,200}$. The strict inequality $A_{200}<1.35$ turns the
limiting bound into an eventual bound. Indeed, let
$\varepsilon=(1.35-A_{200})/2>0$. By the definition of the upper limit,
\[
 \frac{\log m_K}{K^2}\leq A_{200}+\varepsilon<1.35
\]
for all sufficiently large admissible $K$. Consequently we may use
\begin{equation}\label{eq:arith-rounded}
 \boxed{m_KF_K\in\Z[X],\qquad \log m_K\leq1.35K^2.}
\end{equation}

\newpage
\heading{8. The analytic estimate and its integral meaning}{Lai, section 11, Lemma 11.6 and Proposition 11.7.}
The second substantial input concerns the same normalized determinant
$F_K=S_K\Delta_K$.

\textbf{Analytic input.} For every positive multiple $K$ of $40$,
\begin{equation}\label{eq:analytic-input}
 \log F_K(\xi)\leq UK^2+24K\log K+200K,
 \qquad U=-\frac{2733991}{2000000}=-1.3669955.
\end{equation}
Positivity of $F_K(\xi)$ was proved above, so its logarithm is defined.
This is Lai's Proposition 11.7, taken as given with its required uniform
energy bounds. We are not replacing those bounds by numerical sampling.

Dividing by $K^2$ gives
\[
 \frac{\log F_K(\xi)}{K^2}
 \leq-1.3669955+\frac{24\log K}{K}+\frac{200}{K}.
\]
The last two terms tend to zero. Since
$1.3669955-1.36=0.0069955>0$, they are eventually smaller than this gap.
Exponentiating yields
\begin{equation}\label{eq:analytic-rounded}
 \boxed{0<F_K(\xi)\leq\exp(-1.36K^2)}
\end{equation}
for all sufficiently large $K$. No explicit threshold is needed for an
irrationality proof.

To explain what is being estimated, we also state the appropriate form of
Andr\'eief's identity from section 11. For the Gram matrix already constructed,
\begin{equation}\label{eq:andreief}
 \Delta_K(\xi)=\frac1{h!}\int_{(0,\infty)^h}
 \prod_{i<j}(y_i^2-y_j^2)^2
 \prod_{i=1}^h A_K(y_i^2)w(y_i)\dd y_1\cdots\dd y_h.
\end{equation}
This is an $h$-fold integral, not a determinant moved inside the one-variable
matrix integral. The pairwise product is the square of the Vandermonde
determinant $\det[y_j^{2i}]$, with $0\leq i<h$ and $1\leq j\leq h$.
It is nonnegative and vanishes when two integration points coincide.

The scale $K^2$ becomes natural in this expression. Away from coincident
points, taking the logarithm of the integrand gives
\[
 2\sum_{i<j}\log|y_i^2-y_j^2|
 +\sum_i\log A_K(y_i^2)+\sum_i\log w(y_i).
\]
There are $h(h-1)/2$ pairwise terms, with $h$ proportional to $K$.
After rescaling the variables, these interactions compete with the decay
of the weights. Section 11 controls that balance by comparison with a
carefully chosen measure, then combines the bound with $S_K$.
Equation \eqref{eq:analytic-input} is the quantitative result of that work.

The arithmetic and analytic estimates must concern the same $F_K$.
Changing its scalar normalization changes both costs. Their sum, which
bounds the final integer polynomial, is what matters.

\newpage
\heading{9. Combining the estimates and finishing the proof}{The final implication, corresponding to Lai's section 12.}
For all sufficiently large positive integers $n$, define
\[
 Q_n(X)=m_{40n}F_{40n}(X).
\]
The arithmetic input gives integer coefficients. Both scalar multipliers
$S_{40n}$ and $m_{40n}$ are positive and independent of $X$, so
\[
 Q_n\in\Z[X],\qquad \deg Q_n\leq37n,
 \qquad Q_n(\xi)>0.
\]
Take $n$ large enough that both eventual bounds
\eqref{eq:arith-rounded} and \eqref{eq:analytic-rounded} hold at $K=40n$.
Then
\begin{align*}
 \log Q_n(\xi)
 &=\log m_K+\log F_K(\xi)\\
 &\leq1.35K^2-1.36K^2
 =-\frac{K^2}{100}=-16n^2.
\end{align*}
Thus $0<Q_n(\xi)\leq e^{-16n^2}$. The opening criterion applies with
$c=16$. More explicitly, if $\xi=a/b$, then the positive integer
$b^{37n}Q_n(a/b)$ satisfies
\[
 0<b^{37n}Q_n(a/b)
 \leq\exp(37n\log b-16n^2).
\]
For $n\geq37\log b/8$, the exponent is at most $-8n^2$, and the upper
bound is smaller than one. This is the contradiction. Subject to the two
stated quantitative inputs, $\zeta(5)$ is irrational.

The determinant's degree grows linearly, while the small-value estimate
has a quadratic exponent. Positivity prevents vanishing, and the arithmetic
bound leaves enough of that decay after the coefficients become integral.
The particular constant $16$ is only a convenient consequence of rounding;
any fixed positive decay constant would finish the proof.

\bigskip
\textbf{Source note and scope of the conclusion.}
Lai's Theorem 9.3 and equations (106)--(107) include an additional term
$+E_M$ in the arithmetic cost. His following prose instead treats that
positive term as a subtraction. Keeping the printed formulation would give
\[
 \limsup_{n\to\infty}\frac{\log Q_n(\xi)}{(40n)^2}
 \leq A_{200}+U+E_{200}
\]
for the polynomials of that formulation. To infer irrationality from that
bound requires the additional inequality
\[
 E_{200}<-A_{200}-U=0.017407802258295\ldots.
\]
It would yield some positive decay constant, not necessarily $16$.
This account assumes Fauzan's original arithmetic bound explicitly.
That source choice does not verify or repair the disputed arithmetic, and
the conclusion is not claimed to follow solely from Lai's literal corrected
sections 8--11.

\source{Sources: Chun Wing Lai,
\href{https://ai.vixra.org/pdf/2609.0080v1.pdf}{\emph{On Fauzan's Proof that $\zeta(5)$ is Irrational}},
v1, 29 September 2026; Aabir Fauzan,
\href{https://zenodo.org/records/22826419}{\emph{$\zeta(5)$ is irrational}},
17 September 2026, especially Proposition 5.1, (5.21), and Appendix B.
The sections of this account have their own numbering.}

\newpage
\heading{Appendix A. A common source for the two integral formulas}
{A Fourier-series calculation supplies the classical identities used in the proof.}

The moment and pole calculations use two familiar facts: Euler's formula
for even zeta values and a special Hurwitz zeta integral. Both follow from
one elementary identity. Put $f(y)=(e^{2\pi y}-1)^{-1}$. For $u>0$,
\begin{equation}\label{eq:kernel}
 \frac1{1-e^{-u}}=\frac12+\frac1u
       +2\int_0^\infty f(y)\sin(uy)\dd y.
\end{equation}

To prove it, periodically extend
$F_u(x)=\cosh(u(x-\tfrac12))$ from $[0,1]$. Its endpoints agree,
and its Fourier coefficients satisfy
\[
 c_n=\int_0^1 F_u(x)e^{-2\pi\mathrm i nx}\dd x
     =\frac{2u\sinh(u/2)}{u^2+4\pi^2n^2},\qquad n\in\Z.
\]
Indeed, integrate $F_u''=u^2F_u$ twice by parts. The equal endpoint
values cancel, and the endpoint derivative difference is
$2u\sinh(u/2)$, giving
$(u^2+4\pi^2n^2)c_n=2u\sinh(u/2)$.
The periodic function is continuous and piecewise smooth. The Fourier
convergence theorem at $x=0$ therefore gives
$\cosh(u/2)=c_0+2\sum_{n\geq1}c_n$. Dividing by $2\sinh(u/2)$ yields
\[
 \frac12\coth(u/2)=\frac1u
      +2u\sum_{n\geq1}\frac1{u^2+4\pi^2n^2}.
\]
Expanding $f$ as a geometric series and integrating each exponential gives
\[
 \int_0^\infty f(y)\sin(uy)\dd y
 =\sum_{n\geq1}\int_0^\infty e^{-2\pi ny}\sin(uy)\dd y
 =\sum_{n\geq1}\frac{u}{u^2+4\pi^2n^2}.
\]
This interchange is justified by $|\sin(uy)|\leq uy$:
the sum of the absolute integrals is at most
$u\sum_{n\geq1}(2\pi n)^{-2}<\infty$.
Together with $1/(1-e^{-u})=\tfrac12+\tfrac12\coth(u/2)$,
this proves \eqref{eq:kernel}.

\textbf{Euler's formula.}
The same partial fractions give, for $0<u<2\pi$,
\begin{align*}
 \frac{u}{e^u-1}
 &=1-\frac u2+2u^2\sum_{n\geq1}\frac1{u^2+4\pi^2n^2}\\
 &=1-\frac u2+
   2\sum_{r\geq0}\frac{(-1)^r\zeta(2r+2)}{(2\pi)^{2r+2}}u^{2r+2}.
\end{align*}
For the second equality, expand each denominator geometrically. The
resulting double series converges absolutely whenever $|u|<2\pi$, so
the order of summation may be exchanged. Compare with the defining
series $u/(e^u-1)=\sum_{k\geq0}B_ku^k/k!$. Equality of coefficients gives
\[
 \boxed{\zeta(2e+2)=
 \frac{(-1)^eB_{2e+2}(2\pi)^{2e+2}}{2(2e+2)!}},
 \qquad e\in\Z_{\geq0}.
\]
This is the identity used to turn the polynomial integrals into the
rational Bernoulli values defining $\mu$.

\newpage
\heading{Appendix A, continued. The Hurwitz integral}
{Deriving the fifth-order formula, including the interchange of integrals.}

Fix $a>0$. The definition
$\zeta(5,a)=\sum_{n\geq0}(a+n)^{-5}$ and the elementary exponential
integral give
\[
 \frac1{(a+n)^5}=\frac1{24}
       \int_0^\infty u^4e^{-(a+n)u}\dd u.
\]
The integrands are nonnegative, so Tonelli's theorem allows the sum and
integral to be exchanged. Summing the geometric series gives
\[
 \zeta(5,a)=\frac1{24}
       \int_0^\infty\frac{u^4e^{-au}}{1-e^{-u}}\dd u.
\]
Insert \eqref{eq:kernel}. The first two terms can be evaluated immediately:
\[
 \frac1{24}\int_0^\infty u^4e^{-au}
       \left(\frac12+\frac1u\right)\dd u
       =\frac1{2a^5}+\frac1{4a^4}.
\]
The remaining term is a double integral. Although the sine can change
sign, Fubini's theorem applies because
\begin{align*}
 &\int_0^\infty\!\int_0^\infty
       u^4e^{-au}f(y)|\sin(uy)|\dd y\dd u\\
 &\hspace{12mm}\leq
       \left(\int_0^\infty u^5e^{-au}\dd u\right)
       \left(\int_0^\infty yf(y)\dd y\right)<\infty.
\end{align*}
The first factor is $5!/a^6$. For the second, $yf(y)$ tends to
$1/(2\pi)$ at zero and decays exponentially at infinity. Thus the bound
also justifies subtracting and rearranging the integrals in the preceding
calculation.

After exchanging the integrations, evaluate the inner integral by taking
the imaginary part of an exponential integral:
\[
 \int_0^\infty u^4e^{-au}\sin(uy)\dd u
 =\operatorname{Im}\int_0^\infty
       u^4e^{-(a-\mathrm i y)u}\dd u
 =24\operatorname{Im}(a-\mathrm i y)^{-5}.
\]
The last equality follows by four integrations by parts; every boundary
term vanishes because the real part of $a-\mathrm i y$ is positive.
Keeping the factor $2/24$ from \eqref{eq:kernel}, we obtain
\[
 \boxed{\zeta(5,a)=\frac1{2a^5}+\frac1{4a^4}
 +2\int_0^\infty
   \frac{\operatorname{Im}(a-\mathrm i y)^{-5}}{e^{2\pi y}-1}\dd y.}
\]
This is precisely the special Hermite formula \eqref{eq:hermite} used in
the pole calculation. In particular, it holds for every positive integer
$a=j$, where the Hurwitz zeta value is the tail of the series for
$\zeta(5)$.

The two classical identities have now been derived from the same kernel.
The remaining passage to the positive weight $w$ is the integration by
parts carried out in the main text. Together, these calculations justify
the integral representation on both parts of the partial-fraction basis.

\bigskip
{\small\textbf{Classical references.}
The identities derived here are also recorded in the NIST Digital Library
of Mathematical Functions:
\href{https://dlmf.nist.gov/4.36.E3}{the partial fractions for $\coth$ (4.36.3)},
\href{https://dlmf.nist.gov/25.6.E2}{Euler's even-zeta formula (25.6.2)}, and
\href{https://dlmf.nist.gov/25.11.E29}{Hermite's Hurwitz integral (25.11.29)}.}

\newpage
\heading{Appendix B. Why the degree is exactly $h$}
{Lai, section 6, Proposition 6.1. Optional for the irrationality argument.}

The main proof only needs $\deg\Delta_K\leq h$. We can also show that
the highest-degree term does not cancel. The reason is that cancellation
in $A_K=D_N^6/D_K$ leaves exactly $h=K-N$ distinct poles, matching the
number of rows of the matrix.

Write $a_k=-k^2$ for $N<k\leq K$ and
$T_K(t)=\prod_{k=N+1}^K(t-a_k)$. The residue of
$A_K(t)t^{i+j}=D_N(t)^5t^{i+j}/T_K(t)$ at $a_k$ is
\[
 c_k^{ij}=\frac{D_N(a_k)^5a_k^{i+j}}{T_K'(a_k)}.
\]
Indeed, multiply by $t-a_k$ and evaluate at $a_k$; the other denominator
factors multiply to $T_K'(a_k)$. The polynomial part of the partial
fraction decomposition contributes no $X$ to $\mu_X$. By
\eqref{eq:mu-pole}, the coefficient of $X$ contributed by this pole
is $k^4c_k^{ij}$. Thus, writing $G_K(X)=G_0+XG_1$, we have
\[
 (G_1)_{ij}=\sum_{k=N+1}^K d_ka_k^ia_k^j,
 \qquad d_k=\frac{k^4D_N(a_k)^5}{T_K'(a_k)}.
\]

Let $V$ be the square matrix with entries $V_{ik}=a_k^i$, where
$0\leq i<h$ and $N<k\leq K$. The preceding equality is the matrix identity
\[
 G_1=V\operatorname{diag}(d_k)V^{\mathsf T},\qquad
 \det V=\prod_{N<k<\ell\leq K}(a_\ell-a_k).
\]
For completeness, the Vandermonde formula follows because the determinant
is an alternating polynomial in the nodes: it vanishes whenever two nodes
coincide, so every difference $a_\ell-a_k$ divides it. The product of these
differences has the same total degree $0+1+\cdots+(h-1)$. Comparing the
coefficient of the monomial whose successive exponents are
$0,1,\ldots,h-1$ shows that the remaining constant is one.

All nodes are distinct, so $\det V\ne0$. Also $d_k\ne0$: the factors
of $D_N(a_k)=\prod_{j=1}^N(j^2-k^2)$ are nonzero, as are those of
$T_K'(a_k)$. Therefore $\det G_1\ne0$. In the determinant expansion,
the coefficient of $X^h$ is obtained by taking the $X$ term from each of
the $h$ factors. Consequently
\[
 [X^h]\Delta_K(X)=\det G_1\ne0,\qquad \deg\Delta_K=h.
\]

The leading coefficient has an especially simple form. Pairing the two
factors belonging to each unordered pair of nodes gives
\[
 \prod_{k=N+1}^K T_K'(a_k)
 =(-1)^{h(h-1)/2}\prod_{N<k<\ell\leq K}(a_\ell-a_k)^2.
\]
Using this in $\det G_1=(\det V)^2\prod_k d_k$ cancels the squared
differences and yields
\[
 \boxed{[X^h]\Delta_K(X)
 =(-1)^{h(h-1)/2}\prod_{k=N+1}^K k^4D_N(-k^2)^5
 \in\Z\setminus\{0\}.}
\]
This concerns one coefficient. It does not assert that all coefficients
of $\Delta_K$ are integers, and it is separate from the positivity of
$\Delta_K(\xi)$.

\newpage
\heading{Appendix C. The same functional in different coordinates}
{Lai, section 7, Lemma 7.3. Optional once the arithmetic estimate is assumed.}

The auxiliary functional in section 7 changes the pole locations from
negative squares to integers. Here is the complete identity, including
the weight that accompanies the substitution.

Define $L:\Q[x]\to\Q$ by $L(x^k)=B_k$, and put
\[
 \tau(P)=\frac{L(P''')}{24},\qquad
 d(r)=\begin{cases}r,&r\geq0,\\-r-1,&r<0.\end{cases}
\]
On rational functions with only simple integer poles, define $\tau_X$
by its polynomial restriction $\tau$ and the pole values
\[
 \tau_X\!\left(\frac1{x-r}\right)=\Hfive{d(r)}-X,
 \qquad r\in\Z.
\]
Unique partial fractions extend these prescriptions linearly, just as
in the definition of $\mu_X$.

For every rational function $R$ in the domain of $\mu_X$, the identity is
\[
 \boxed{\mu_X(R)=\tau_X\!\left(x^5R(-x^2)\right).}
\]
The right side is defined: a pole at $t=-j^2$ becomes a pair of simple
poles at $x=j$ and $x=-j$, and a polynomial remains a polynomial after
substitution and multiplication by $x^5$. Both sides are linear in $R$.
By \eqref{eq:partial}, it is enough to check monomials and simple poles.

For $R(t)=t^e$, the transformed function is $(-1)^ex^{2e+5}$. Three
differentiations give
\[
 \tau_X\bigl((-1)^ex^{2e+5}\bigr)
 =\frac{(-1)^e(2e+5)(2e+4)(2e+3)B_{2e+2}}{24},
\]
which is exactly \eqref{eq:mu-moment}. This explains the order of the
derivative: it produces both the three consecutive factors and the
required Bernoulli index.

For $R(t)=1/(t+j^2)$, polynomial division gives
\[
 \frac{x^5}{j^2-x^2}
 =-x^3-j^2x-\frac{j^4}{2}
       \left(\frac1{x-j}+\frac1{x+j}\right).
\]
One can check the signs by multiplying by $j^2-x^2$. The polynomial
terms give $x^5-j^4x$, and the pole terms give $+j^4x$.
Now $\tau(-x^3)=-1/4$, $\tau(-j^2x)=0$, $d(j)=j$, and $d(-j)=j-1$.
Applying $\tau_X$ therefore gives
\begin{align*}
 \tau_X\!\left(\frac{x^5}{j^2-x^2}\right)
 &=-\frac14-\frac{j^4}{2}
       \left(\Hfive{j}+\Hfive{j-1}-2X\right)\\
 &=j^4\bigl(X-\Hfive{j}\bigr)-\frac14+\frac1{2j},
\end{align*}
where the last equality uses
$\Hfive{j}-\Hfive{j-1}=j^{-5}$. This is \eqref{eq:mu-pole}, completing
the proof by linearity.

Both the substitution $t=-x^2$ and the factor $x^5$ are necessary.
The identity supplies a second description of the same functional;
it introduces no new determinant. Integer pole locations make the
arithmetic calculations by residue class possible, while the original
$t$-description makes the positive integral representation visible.

\end{document}
